\chapter{2023年上海卷（春）}

\section{填空题}

\begin{problem}
已知集合 \(A=\{1,2\}\)，\(B=\{1,a\}\)，且 \(A=B\)，则 \(a=\)\fillinblank{} .
\end{problem}

\begin{problem}
已知向量 \(\overrightarrow{a}=(3,4)\)，\(\overrightarrow{b}=(1,2)\)，则 \(\overrightarrow{a}-2\overrightarrow{b}=\)\fillinblank{} .
\end{problem}

\begin{problem}
若不等式 \(\left|x-1\right|\le 2\)，则实数 \(x\) 的取值范围为\fillinblank{} .
\end{problem}

\begin{problem}
已知圆 \(C\) 的一般方程为 \(x^2+2x+y^2=0\)，则圆 \(C\) 的半径为\fillinblank{} .
\end{problem}

\begin{problem}
已知事件 \(A\) 发生的概率为 \(P(A)=0.5\)，则它的对立事件 \(\bar{A}\) 发生的概率 \(P(\bar{A})=\)\fillinblank{} .
\end{problem}

\begin{problem}
已知正实数 \(a\)，\(b\) 满足 \(a+4b=1\)，则 \(ab\) 的最大值为\fillinblank{} .
\end{problem}

\begin{problem}
某校抽取 \(100\) 名学生测身高，其中身高最大值为 \(186\) cm，最小值为 \(154\) cm，根据身高数据绘制频率组距分布直方图，组距为 \(5\)，且第一组下限为 \(153.5\)，则组数为\fillinblank{} .
\end{problem}

\begin{problem}
设 \(\left(1-2x\right)^4=a_0+a_1x+a_2x^2+a_3x^3+a_4x^4\)，则 \(a_0+a_4=\)\fillinblank{} .
\end{problem}

\begin{problem}
已知函数 \(f(x)=2^{-x}+1\)，且 \(g(x)= \begin{cases}
\log_2(x+1),  &x\ge 0,\\
f(-x),  &x<0,
\end{cases}\)
则方程 \(g(x)=2\) 的解为\fillinblank{} .
\end{problem}

\begin{problem}
已知有 \(4\) 名男生 \(6\) 名女生，现从 \(10\) 人中任选 \(3\) 人，则恰有 \(1\) 名男生 \(2\) 名女生的概率为\fillinblank{} .
\end{problem}

\begin{problem}
设 \(z_1\)，\(z_2\in \mathbb{C}\) 且 \(z_1=\i\cdot z_2\)，满足 \(\left|z_1-1\right|=1\)，则 \(\left|z_1-z_2\right|\) 的取值范围为\fillinblank{} .
\end{problem}

\begin{problem}
已知空间向量 \(\overrightarrow{OA}\)，\(\overrightarrow{OB}\)，\(\overrightarrow{OC}\) 都是单位向量，且 \(\overrightarrow{OA}\perp \overrightarrow{OB}\)，\(\overrightarrow{OA}\perp \overrightarrow{OC}\)，\(\overrightarrow{OB}\) 与 \(\overrightarrow{OC}\) 的夹角为 \(60^\circ\). 若 \(P\) 为空间任意一点，且 \(\left|\overrightarrow{OP}\right|=1\)，满足 \(\left|\overrightarrow{OP}\cdot\overrightarrow{OC}\right|\le \left|\overrightarrow{OP}\cdot\overrightarrow{OB}\right|\le \left|\overrightarrow{OP}\cdot\overrightarrow{OA}\right|\)，则 \(\overrightarrow{OP}\cdot\overrightarrow{OC}\) 的最大值为\fillinblank{} .
\end{problem}

\section{单选题}

\begin{problem}
下列函数是偶函数的是（ \quad ）.

\choices
    {\(y=\sin x\)}
    {\(y=\cos x\)}
    {\(y=x^3\)}
    {\(y=2^x\)}
\end{problem}

\begin{problem}
如图为 \(2018-2021\) 年中国货物进出口总额的条形统计图，则下列对于进出口贸易额描述错误的是（ \quad ）.

\bitmapfigure[max width=0.48\linewidth]{img/2023/shanghai_spring/q14_fig1.png}

\choices
    {从 \(2018\) 年开始，\(2021\) 年的进出口总额增长率最大}
    {从 \(2018\) 年开始，进出口总额逐年增大}
    {从 \(2018\) 年开始，进口总额逐年增大}
    {从 \(2018\) 年开始，\(2020\) 年的进出口总额增长率最小}
\end{problem}

\begin{problem}
如图所示，在正方体 \(ABCD-A_1B_1C_1D_1\) 中，点 \(P\) 为线段 \(A_1C_1\) 上的动点，则下列直线中，始终与直线 \(BP\) 异面的是（ \quad ）.

\bitmapfigure[max width=0.38\linewidth]{img/2023/shanghai_spring/q15_fig1.png}

\choices
    {\(DD_1\)}
    {\(AC\)}
    {\(AD_1\)}
    {\(B_1C\)}
\end{problem}

\begin{problem}
已知数列 \(\{a_n\}\) 的各项均为实数，\(S_n\) 为其前 \(n\) 项和，若对任意 \(k>2022\)，都有 \(\left|S_k\right|>\left|S_{k+1}\right|\)，则下列说法正确的是（ \quad ）.

\choices
    {\(a_1,a_3,a_5,\cdots,a_{2n-1}\) 为等差数列，\(a_2,a_4,a_6,\cdots,a_{2n}\) 为等比数列}
    {\(a_1\)，\(a_3\)，\(a_5\)，\(\cdots\)，\(a_{2n-1}\) 为等比数列，\(a_2\)，\(a_4\)，\(a_6\)，\(\cdots\)，\(a_{2n}\) 为等差数列}
    {\(a_1\)，\(a_2\)，\(a_3\)，\(\cdots\)，\(a_{2022}\) 为等差数列，\(a_{2022}\)，\(a_{2023}\)，\(a_{2024}\)，\(\cdots\)，\(a_n\) 为等比数列}
    {\(a_1\)，\(a_2\)，\(a_3\)，\(\cdots\)，\(a_{2022}\) 为等比数列，\(a_{2022}\)，\(a_{2023}\)，\(a_{2024}\)，\(\cdots\)，\(a_n\) 为等差数列}
\end{problem}

\section{解答题}

\begin{problem}
已知三棱锥 \(P-ABC\) 中，\(AB\perp AC\)，\(PA\perp\) 平面 \(ABC\)，\(PA=AB=3\)，\(AC=4\)，\(M\) 为 \(BC\) 中点，过点 \(M\) 分别作平行于平面 \(PAB\) 的直线交 \(AC\)，\(PC\) 于点 \(E\)，\(F\).

\begin{enumerate}
\item 求直线 \(PM\) 与平面 \(ABC\) 所成的角的正切值；
\item 证明：平面 \(MEF\parallel\) 平面 \(PAB\)，并求直线 \(ME\) 到平面 \(PAB\) 的距离.
\end{enumerate}

\FigureLayoutDeclare{img/2023/shanghai_spring/q17_fig1.png}{5.0cm}{20bp 31bp 60bp 5bp}
\bitmapfigure[max width=0.42\linewidth]{img/2023/shanghai_spring/q17_fig1.png}
\end{problem}

\begin{problem}
在 \(\triangle ABC\) 中，内角 \(A\)，\(B\)，\(C\) 所对的边为 \(a\)，\(b\)，\(c\)，其中 \(b=2\).

\begin{enumerate}
\item 若 \(A+C=60^\circ\)，且 \(a=2c\)，求边长 \(c\) 的值；
\item 若 \(A-C=30^\circ\)，\(a=\sqrt2\,c\sin A\)，求 \(S_{\triangle ABC}\).
\end{enumerate}
\end{problem}

\begin{problem}
已知 \(S\) 为正比例系数，定义 \(S=\frac{F_0}{V_0}\)，\(F_0\) 为建筑物暴露在空气中的面积（单位：平方米），\(V_0\) 为建筑物的体积（单位：立方米）.

\begin{enumerate}
\item 若有一个圆柱体建筑的底面半径为 \(R\)，高度为 \(H\)，求该建筑体的 \(S\)（用 \(R\)，\(H\) 表示）；
\item 现有一个建筑体，侧面皆垂直于地面，设 A 为底面面积，\(L\) 为建筑底面周长. 已知 \(f\) 为正比例系数，\(L^2\) 与 A 成正比，定义 \(f=\frac{L^2}{A}\)，建筑面积即为每一层的底面面积，总建筑面积即为每层建筑面积之和，值为 \(T\). 已知该建筑体推导得出 \(S=\sqrt{\frac{f\cdot n}{T}+\frac{1}{3n}}\)，\(n\) 为层数，层高为 \(3\) 米，其中 \(f=18\)，\(T=10000\)，试求当取第几层时，该建筑体 \(S\) 最小？
\end{enumerate}
\end{problem}

\begin{problem}
已知椭圆 \(\Gamma:\frac{x^2}{m^2}+\frac{y^2}{3}=1\) \(\left(m>0,\ m\neq \sqrt3\right)\).

\begin{enumerate}
\item 若 \(m=2\)，求椭圆 \(\Gamma\) 的离心率；
\item 设 \(A_1\)，\(A_2\) 为椭圆 \(\Gamma\) 的左右顶点，若椭圆 \(\Gamma\) 上一点 \(E\) 的纵坐标为 \(1\)，且 \(\overrightarrow{EA_1}\cdot\overrightarrow{EA_2}=-2\)，求 \(m\) 的值；
\item 若 \(P\) 为椭圆 \(\Gamma\) 上一点，过点 \(P\) 作一条斜率为 \(\sqrt3\) 的直线与双曲线 \(\frac{y^2}{5m^2}-\frac{x^2}{5}=1\) 仅有一个公共点，求 \(m\) 的取值范围.
\end{enumerate}
\end{problem}

\begin{problem}
设函数 \(f(x)=ax^3-(a+1)x^2+x\)，\(g(x)=kx+m\)，其中 \(a\ge 0\)，\(k\)，\(m\in \R\)，若任意 \(x\in[0,1]\) 均有 \(f(x)\le g(x)\)，则称函数 \(y=g(x)\) 是函数 \(y=f(x)\) 的控制函数，且对于所有满足条件的函数 \(y=g(x)\) 在 \(x\) 处取得的最小值记为 \(\bar f(x)\).

\begin{enumerate}
\item 若 \(a=2\)，\(g(x)=x\)，试问 \(y=g(x)\) 是否为 \(y=f(x)\) 的控制函数；
\item 若 \(a=0\)，使得直线 \(y=h(x)\) 是曲线 \(y=f(x)\) 在 \(x=\frac14\) 处的切线，证明：函数 \(y=h(x)\) 为函数 \(y=f(x)\) 的控制函数，并求 \(\bar f\left(\frac14\right)\) 的值；
\item 若曲线 \(y=f(x)\) 在 \(x=x_0\)（\(x_0\in(0,1)\)）处的切线过点 \((1,0)\)，且 \(c\in[x_0,1]\)，证明：当且仅当 \(c=x_0\) 或 \(c=1\) 时，\(\bar f(c)=f(c)\).
\end{enumerate}
\end{problem}

